\clearpage
\fancyhead[L]{\small INTRODUCTION}
\SourceLeaf{351}
\pdfbookmark[0]{Introduction to the Magnetic Association's 1836 Observations}{einleitung-magnetic-association-1836}
\begin{center}
{\Large INTRODUCTION}\par\medskip
{\small [FOR THE JOURNAL: RESULTS OF THE OBSERVATIONS OF THE MAGNETIC ASSOCIATION IN 1836. EDITED BY C. F. GAUSS AND W. WEBER.]}
\end{center}
\medskip
Of the many manifestations of the Earth's magnetic force, which can only be understood through numerous precise observations carried on at widely different points on the Earth's surface, the \textit{irregular changes} that we find in this force call most for strictly coordinated work among observers. It is well known that the elements that determine the Earth's magnetic force, namely declination, inclination and no doubt also intensity, undergo continual changes. Adequate experience is still lacking for intensity, which has only been studied for a few decades. The changes include secular ones, which become noticeable only after a longer interval but grow quite considerable with time, and periodic ones, which vary with the time of year and day. Insofar as these changes occur regularly, strictly coordinated work by observers at different places is highly desirable for advancing our understanding more quickly, but it is unnecessary in principle. Each observer can also make useful contributions independently of the others.

The situation is different for irregular changes, which have begun to receive greater attention only in recent years. Al\SourceLeaf{352}most a hundred years ago, \textsc{Hiorter} and \textsc{Celsius} noticed that when an aurora is visible the magnetic needle moves irregularly, often by a very large amount. Many observations subsequently confirmed this. One could conclude that the same forces that produce the appearance of an aurora also act on the magnetic needle, and that these effects must extend over very great distances, since the northern lights are usually visible over a wide area. Mr. \textsc{Arago}'s observation gives us an even better idea of how far those puzzling forces reach. Often, on the very days when he had observed severe disturbances of the regular motion of the magnetic needle in Paris, northern lights had been seen in distant places without becoming visible above the Paris horizon.

This gave a particular interest to the irregularities in the manifestations of terrestrial magnetism. Mr. \textsc{von Humboldt}, in particular, had noticed how often they occurred in his many observations of the daily and hourly motions of the magnetic needle. Those remarks gave no grounds for treating all irregular movements as simultaneous with auroras and left open the possibility that many, perhaps most, were due solely to local causes. Yet it was hard to miss that powerful natural forces acting over long distances are not infrequently at work. Understanding how those forces act and how far they reach is a worthy subject of research into nature, even if for the time being their source remains unknown.

Observations made in passing and left to chance cannot bring us closer to this goal. To reach it, many such phenomena must be followed simultaneously and in exact detail at many places, with precise measurements of their timing and magnitude. This requires specific advance arrangements among observers who have suitable instruments at their disposal.

The famous naturalist to whom we owe so much of what we know about terrestrial magnetism was also the first to break new ground here. Toward the end of 1828, Mr. \textsc{von Humboldt} built a small iron-free building in Berlin specifically for magnetic observations. He installed a variation compass made by \textsc{Gambey} and joined with other own\SourceLeaf{353}ers of similar instruments in several, in some cases very distant, places to observe magnetic variation regularly on agreed days. Eight observing sessions a year were set, each lasting 44 hours. During them, magnetic declination was to be recorded every hour. At some places observations were made at even shorter intervals, every half hour or every twenty minutes. The details can be found on page 361 of volume 19 of \textsc{Poggendorff}'s \textit{Annalen der Physik}, along with the observations made under this agreement at the sessions in 1829 and 1830 in Berlin, Freiberg, Petersburg, Kasan and Nicolajef, and graphs of three of those sets of observations.

The magnetic observatory here was built in 1833, and its magnetic apparatus differs entirely in design from what had been used before. The first full set of observations at scheduled sessions here was made on 20 and 21 March 1834. Only Berlin was known to have made corresponding observations. In Göttingen, however, observations were taken every ten minutes, and in Berlin only every hour. Even so, the Berlin records showed several fairly large movements that were also found in the Göttingen observations. The latter showed a great many other movements in the intervals, which of course could leave no trace in the Berlin records. It therefore remained undecided whether a smaller or a larger share of the fluctuations observed in Göttingen had been purely local.

The very next session, on 4 and 5 May, settled the matter. The intervals were shortened further to five minutes, which made the results considerably clearer. No corresponding observations with \textsc{Gambey}'s instruments from this session or any later one have come to light at all. Mr. \textsc{Sartorius}, however, had taken an active part in the March observations at Göttingen and had obtained an instrument just like the Göttingen one, only smaller, for a journey to Italy that was to last several years. With it he ob\SourceLeaf{354}served the May session at short intervals, carefully and in full, at Waltershausen (an estate in Bavaria about 20 Meilen from Göttingen). The agreement was truly surprising. It extended to the larger fluctuations and to almost all the smaller fluctuations that varied over short intervals, leaving, in fact, nothing that one would have been justified in attributing to local causes.

The next three sessions, in June, August and September, were conducted in Göttingen in exactly the same way, while the number of participants elsewhere with similar or identical instruments kept growing. After Professor \textsc{Encke} had seen the arrangements here for himself, he had a similar instrument made on a smaller scale for provisional use in Berlin. Mr. \textsc{Sartorius} observed with his instrument at every session where circumstances allowed (in June in Frankfurt, in September in Bramberg in the Salzburg region). In Leipzig, Copenhagen and Braunschweig observations were made with instruments identical to the one here. The corresponding observations gave a result very similar to that reported for the May session. Almost all of the many fluctuations seen in Göttingen also appeared in the observations elsewhere, in differing relative magnitudes, yet in unmistakable agreement.

To obtain further indisputable evidence of this remarkable result, some special corresponding observations were arranged between Göttingen and Leipzig while Professor \textsc{Weber} was in Leipzig. Specific hours in the morning, at noon and in the evening of 1 and 2 October were set aside for them. These observations were made by especially experienced observers with the greatest care. They are printed in full in \textsc{Poggendorff}'s \textit{Annalen der Physik}, vol.~33, p.~426, and illustrated with graphs.

This showed very clearly the need to follow these phenomena at much shorter intervals than Mr. \textsc{von Humboldt} had chosen. For a time we conducted the scheduled observations at intervals of three minutes, as did some other participants. However, some participants elsewhere kept to five-minute intervals, which also suffice in most ordinary cases, so we later adopted those intervals as the general rule for the sake of uniformity. But at such short intervals, conducting the \SourceLeaf{355}sessions becomes incomparably more demanding than making hourly records, especially where only a few people have to share the work. It was therefore necessary. in order to keep the association going, to reduce both the number and the duration of the sessions. Since then their number has been set at six a year and the duration of each at 24 hours. Two supplementary sessions were added to each main session. More details are given below.

Under this arrangement, observations have continued without interruption and still continue in Göttingen and in an increasing number of other places. Instruments identical or similar to the one at Göttingen are found in Altona, Augsburg, Berlin, Bonn, Braunschweig, Breda, Breslau, Cassel, Copenhagen, Dublin, Freiberg, Göttingen, Greenwich, Halle, Kasan, Krakau, Leipzig, Mailand, Marburg, München, Neapel, Petersburg and Upsala. No observations from eight of these places have yet come to our attention, and at some of the others outside circumstances have so far kept participation from becoming regular and uninterrupted.

Graphs of some sessions from the early days of this association have been published in \textsc{Schumacher}'s \textit{Astronomische Nachrichten} and \textsc{Poggendorff}'s \textit{Annalen der Physik}. Since participation had already grown so much, it seemed time to arrange regular publication, so that this wealth of fruitful facts could become common knowledge among members of the public interested in the study of nature. What we present now can be regarded as the first annual volume. since the association reached a certain size. From 1837 onward, the results of each session will be published regularly, as soon as they have been collected in sufficient completeness.

The observations and their graphs will be accompanied by explanations and remarks directly concerning them, as well as by other papers. These will deal with topics from the broad field of terrestrial magnetism, the instruments used for it, their adjustment and handling, and. the many uses to which they can be put.

\SourceLeaf{356}A further remark may be permitted concerning the immediate subject of our association's work, the changes in magnetic declination. If, as there is no reason to doubt, the two other elements of the Earth's magnetic force, inclination and intensity, are subject to similar changes, one may ask why such careful efforts are devoted chiefly, or for the present exclusively, to the first element.

Knowledge of changes and disturbances in magnetic declination is indeed of very great practical interest. It must matter immensely to the navigator, the geodesist and the mine surveyor to know how often and how severely one of their main instruments is inevitably disturbed, if only to gauge how much trust they can place in it. For the last two uses of the compass, in practical geometry above and below ground, these investigations may prove far more useful still. If it is established that the irregular disturbances that vary with time are never, or only extremely rarely, purely local, and instead always or nearly always occur over wide areas at exactly the same time and with almost the same magnitude, there is a way to make them almost entirely harmless. The geodesist and the mine surveyor need only carry out all their compass operations by the clock with exact timing and arrange for simultaneous observations at another place not too far away. Comparing the observations can eliminate those disturbances in the same way that traveling observers make their barometric height measurements independent of irregular changes in barometer readings by comparing them with barometer observations at a fixed place. It goes without saying that the disturbances a compass experiences in iron-bearing mines are outside this discussion.

The reason for giving declination preference over the other elements of terrestrial magnetism is nevertheless to be sought less in these practical considerations than in the current state of the instruments. For the naturalist, the search for laws in natural phenomena has a purpose and a value of its own, and there is a peculiar charm in recognizing measure and harmony in what is appar\SourceLeaf{357}ently quite irregular. In following the wonderful interplay of the constantly shifting changes in declination, the instruments now used leave nothing to be desired in reliability, precision or ease of observation. The means available so far for observing the other two elements fall short of that standard. For now it is therefore too early to include those elements in large-scale investigations. But as soon as the means of observation have been improved to the point where we can detect the changes, especially the rapidly shifting ones, in the other elements of terrestrial magnetism with certainty, follow them easily and measure them precisely, those changes will have the same claim on the combined efforts of naturalists as the changes in declination. One may hope that this time is no longer very far off.

